Electronic sports (e-sports) have grown into a major global domain, with industry reports describing a worldwide audience in the hundreds of millions and a market value above USD 2 billion \autocite{newzoo2024}. In this highly competitive context, the psychological state of \textit{flow} is considered a key performance factor. Within flow theory, flow emerges when perceived player skill and task challenge are well balanced \autocite{inbook,Moneta2012}. This thesis focuses on \textit{League of Legends} (LoL) as its specific e-sports context because LoL combines a large competitive ecosystem with clearly structured ranked tiers, team-based coordination, rapid decision-making, and frequent high-pressure events such as fights and objective contests \autocite{riotgames_lol,gong2020league,10.1145/3549506}. These characteristics make LoL a suitable case for examining how flow can be detected and protected in a demanding real-time game environment.

From a sensing perspective, the thesis places particular emphasis on electrocardiography (ECG). ECG-derived heart rate and heart rate variability (HRV) features provide a comparatively unobtrusive way to capture autonomic arousal and cognitive load during gameplay, and prior work has demonstrated that flow-related states can be classified from HRV-based physiological signals \autocite{9296229}. Because ECG can be recorded with a chest belt without requiring players to interrupt play or wear highly intrusive equipment, it is especially relevant for e-sports settings in which comfort, movement freedom, and competitive integrity are important. In addition to ECG, webcam-based facial recordings provide non-invasive affect-related cues that can be synchronized with gameplay and physiological data in gaming contexts \autocite{kutt2022biraffe2}. This combination motivates a practical sensing setup that links ECG-derived HRV features with webcam-based affect features for later live evaluation. This technical foundation supports the development of \textit{physio-adaptive systems} that use real-time physiological and affect-related data to infer user states and trigger digital environmental adaptations \autocite{irshad2023wearable,ubt_eref88160}.

Despite the demonstrated feasibility of physiological state detection, research regarding specialized interventions in the e-sports domain remains remarkably scarce. As shown by \cite{ubt_eref88160}, prior work has explored flow support through live biofeedback to enhance user awareness and emotion regulation during task execution. Likewise, \cite{9296229} demonstrated that binary flow states (high vs. low) can be classified from HRV using machine learning models. However, the unique demands of competitive LoL gameplay---characterized by extreme physiological arousal, rapid-fire decision-making, team communication, and specialized motor patterns---are not yet sufficiently addressed in these enterprise-focused approaches \autocite{banyai2019psychology,poulus2021determinants,10.1145/3549506}.


To address this gap, this thesis investigates the design requirements and technical implementation of a physio-adaptive system specifically tailored to protect and sustain flow in e-sports. By moving beyond general classification models, this study seeks to identify specialized, context-aware mechanisms that effectively assist e-sports athletes in maintaining peak performance. Accordingly, the study addresses the following research questions:

\begin{itemize}
    \item \textbf{RQ1:} How can ECG and webcam-based data be utilized to accurately distinguish between high and low flow states in a competitive LoL environment? \label{rq:classification}
    \item \textbf{RQ2:} What are the specific design requirements and intervention mechanisms for a physio-adaptive system that effectively assists LoL players in maintaining their flow state? \label{rq:adaptation}
\end{itemize}   

To answer these questions, the thesis follows a multi-phase mixed-methods design. First, a comprehensive baseline of system requirements is derived from the literature and a series of semi-structured interviews triangulated across three distinct cohorts: active competitive League of Legends (LoL) players, game design experts, and physiological computing specialists. These diverse empirical perspectives serve to identify context-aware adaptive measures for the e-sports environment, directly informing the development of a functional working prototype. During model development, the classification pipeline draws on multimodal training data and established physiological flow-classification frameworks \autocite{9296229,irshad2023wearable}. In the final live evaluation, however, the prototype is restricted to the signals available from the deployed setup---namely ECG-derived features and webcam-based features only. Finally, the prototype’s practical utility and the effectiveness of its adaptive interventions are evaluated through follow-up qualitative interviews and functional testing. Throughout the study, qualitative data are analyzed using Grounded Theory Methodology (GTM) \autocite{recker2021researchmethods} to anchor technical design decisions in the lived experiences and competitive requirements of the athletes \autocite{delveholempaecher2024openaxialselective}.